% GENERATED FILE. Edit canonical lessons, then run npm run generate:books.
% canonical-source-hash: fnv1a64:70114444fb490ce8
% canonical-lessons: MR-C407-antar, MR-C407-veg, MR-C407-brek, MR-C407-disadarsak, MR-C407-paypit

\chapter{Distance and Speed}
\label{ch:mr-distance-and-speed}
\hlchaptermodality{407}

\begin{chapteropening}
\textbf{By the end of this chapter:} \emph{I can talk about a distance, speed, brakes, signposts and a long walk.}

Complete the last lesson of chapter 407: I can talk about a distance, speed, brakes, signposts and a long walk.
\end{chapteropening}

\section[\texorpdfstring{\mr{अंतर} (antar)}{अंतर (antar)}]{\mr{अंतर} (antar) — a distance; a gap}

\label{lesson:MR-C407-antar}



\begin{quote}
Before the new one: say the Marathi for a flyover, then the Marathi for to gaze at.
\end{quote}

\subsection*{You'll want to know: \mr{अंतर}}
\textbf{\mr{अंतर}} — \emph{antar} — \textquotedblleft{}a distance; a gap\textquotedblright{}.

\begin{quote}
\mr{घर} \mr{ते} \mr{शाळा} \mr{हे} \mr{अंतर} \mr{किती} \mr{आहे}?
\end{quote}

\begin{quote}
\emph{ghar te śāḷā he antar kitī āhe?}
\end{quote}

\begin{quote}
What is the distance from home to school?
\end{quote}

Neuter.

X \mr{ते} Y, from X to Y: \mr{घर} \mr{ते} \mr{शाळा}.

\begin{cousinweb}
\textbf{\mr{अंतर}} (\emph{antar}) is a tatsama, taken from Sanskrit \emph{antara}, \textquotedblleft{}interior; interval, space between\textquotedblright{} (Monier-Williams).
\end{cousinweb}

\begin{grammarlens}[title={What you've built}]
One more word for getting around.
\end{grammarlens}

\subsection*{Your turn}
\begin{itemize}
\raggedright
  \item \emph{Say it:} \emph{antar}
  \item \emph{Say it:} \emph{antar}, once more
  \item \emph{Say it:} \emph{ghar te śāḷā he antar kitī āhe?} — What is the distance from home to school?
  \item \emph{Recall:} say the Marathi for a flyover, then the Marathi for to gaze at, then say \emph{antar} again
\end{itemize}

\begin{tcolorbox}[breakable,colback=teal!4,colframe=teal!35!black,title={Before you move on}]
What does \mr{अंतर} mean? (\textquotedblleft{}A distance; a gap\textquotedblright{}.) Say it once more.
\end{tcolorbox}
\section[\texorpdfstring{\mr{वेग} (veg)}{वेग (veg)}]{\mr{वेग} (veg) — speed}

\label{lesson:MR-C407-veg}



\begin{quote}
Before the new one: say the Marathi for to gaze at, then the Marathi for a distance.
\end{quote}

\subsection*{You'll want to know: \mr{वेग}}
\textbf{\mr{वेग}} — \emph{veg} — \textquotedblleft{}speed\textquotedblright{}.

\begin{quote}
\mr{घाटात} \mr{बसचा} \mr{वेग} \mr{कमी} \mr{असतो}.
\end{quote}

\begin{quote}
\emph{ghāṭāt bascā veg kamī asto.}
\end{quote}

\begin{quote}
On the ghat road the bus goes at a low speed.
\end{quote}

Masculine.

\mr{वेगाने}, at speed, fast; \mr{वेग} \mr{वाढवणे}, to speed up.

\begin{cousinweb}
\textbf{\mr{वेग}} (\emph{veg}) is a tatsama, taken from Sanskrit \emph{vega}, \textquotedblleft{}violent agitation, rush, speed\textquotedblright{} (Monier-Williams).
\end{cousinweb}

\begin{grammarlens}[title={What you've built}]
One more word for getting around.
\end{grammarlens}

\subsection*{Your turn}
\begin{itemize}
\raggedright
  \item \emph{Say it:} \emph{veg}
  \item \emph{Say it:} \emph{veg}, once more
  \item \emph{Say it:} \emph{ghāṭāt bascā veg kamī asto.} — On the ghat road the bus goes at a low speed.
  \item \emph{Recall:} say the Marathi for to gaze at, then the Marathi for a distance, then say \emph{veg} again
\end{itemize}

\begin{tcolorbox}[breakable,colback=teal!4,colframe=teal!35!black,title={Before you move on}]
What does \mr{वेग} mean? (\textquotedblleft{}Speed\textquotedblright{}.) Say it once more.
\end{tcolorbox}
\section[\texorpdfstring{\mr{ब्रेक} (brek)}{ब्रेक (brek)}]{\mr{ब्रेक} (brek) — a brake}

\label{lesson:MR-C407-brek}



\begin{quote}
Before the new one: say the Marathi for a distance, then the Marathi for speed.
\end{quote}

\subsection*{You'll want to know: \mr{ब्रेक}}
\textbf{\mr{ब्रेक}} — \emph{brek} — \textquotedblleft{}a brake\textquotedblright{}.

\begin{quote}
\mr{चालकाने} \mr{अचानक} \mr{ब्रेक} \mr{दाबला}.
\end{quote}

\begin{quote}
\emph{cālakāne acānak brek dāblā.}
\end{quote}

\begin{quote}
The driver suddenly pressed the brake.
\end{quote}

Masculine.

\mr{ब्रेक} \mr{दाबणे} or \mr{ब्रेक} \mr{लावणे}, to brake.

\begin{cousinweb}
\textbf{\mr{ब्रेक}} (\emph{brek}) is a loan from English \emph{brake}.
\end{cousinweb}

\begin{grammarlens}[title={What you've built}]
One more word for getting around.
\end{grammarlens}

\subsection*{Your turn}
\begin{itemize}
\raggedright
  \item \emph{Say it:} \emph{brek}
  \item \emph{Say it:} \emph{brek}, once more
  \item \emph{Say it:} \emph{cālakāne acānak brek dāblā.} — The driver suddenly pressed the brake.
  \item \emph{Recall:} say the Marathi for a distance, then the Marathi for speed, then say \emph{brek} again
\end{itemize}

\begin{tcolorbox}[breakable,colback=teal!4,colframe=teal!35!black,title={Before you move on}]
What does \mr{ब्रेक} mean? (\textquotedblleft{}A brake\textquotedblright{}.) Say it once more.
\end{tcolorbox}
\section[\texorpdfstring{\mr{दिशादर्शक} (diśādarśak)}{दिशादर्शक (diśādarśak)}]{\mr{दिशादर्शक} (diśādarśak) — a signpost}

\label{lesson:MR-C407-disadarsak}



\begin{quote}
Before the new one: say the Marathi for speed, then the Marathi for a brake.
\end{quote}

\subsection*{You'll want to know: \mr{दिशादर्शक}}
\textbf{\mr{दिशादर्शक}} — \emph{diśādarśak} — \textquotedblleft{}a signpost\textquotedblright{}.

\begin{quote}
\mr{चौकात} \mr{एक} \mr{मोठा} \mr{दिशादर्शक} \mr{आहे}.
\end{quote}

\begin{quote}
\emph{caukāt ek moṭhā diśādarśak āhe.}
\end{quote}

\begin{quote}
There is a big signpost at the junction.
\end{quote}

Masculine.

A formal word: a board that shows the way.

\begin{cousinweb}
\textbf{\mr{दिशादर्शक}} (\emph{diśādarśak}) is built from Sanskrit \emph{diśā}, \textquotedblleft{}direction\textquotedblright{}, and \emph{darśaka}, \textquotedblleft{}showing\textquotedblright{} (Monier-Williams).
\end{cousinweb}

\begin{grammarlens}[title={What you've built}]
One more word for getting around.
\end{grammarlens}

\subsection*{Your turn}
\begin{itemize}
\raggedright
  \item \emph{Say it:} \emph{diśādarśak}
  \item \emph{Say it:} \emph{diśādarśak}, once more
  \item \emph{Say it:} \emph{caukāt ek moṭhā diśādarśak āhe.} — There is a big signpost at the junction.
  \item \emph{Recall:} say the Marathi for speed, then the Marathi for a brake, then say \emph{diśādarśak} again
\end{itemize}

\begin{tcolorbox}[breakable,colback=teal!4,colframe=teal!35!black,title={Before you move on}]
What does \mr{दिशादर्शक} mean? (\textquotedblleft{}A signpost\textquotedblright{}.) Say it once more.
\end{tcolorbox}
\section[\texorpdfstring{\mr{पायपीट} (pāypīṭ)}{पायपीट (pāypīṭ)}]{\mr{पायपीट} (pāypīṭ) — a long trudge; a long walk}

\label{lesson:MR-C407-paypit}



\begin{quote}
Before the new one: say the Marathi for a brake, then the Marathi for a signpost.
\end{quote}

\subsection*{You'll want to know: \mr{पायपीट}}
\textbf{\mr{पायपीट}} — \emph{pāypīṭ} — \textquotedblleft{}a long trudge; a long walk\textquotedblright{}.

\begin{quote}
\mr{बस} \mr{नव्हती} \mr{म्हणून} \mr{आमची} \mr{खूप} \mr{पायपीट} \mr{झाली}.
\end{quote}

\begin{quote}
\emph{bas navhatī mhaṇūn āmcī khūp pāypīṭ jhālī.}
\end{quote}

\begin{quote}
There was no bus, so we had a long trudge on foot.
\end{quote}

Feminine.

\mr{पायपीट} \mr{होणे} or \mr{पायपीट} \mr{करणे}; it sounds tiring.

\begin{cousinweb}
\textbf{\mr{पायपीट}} (\emph{pāypīṭ}) is built from \textbf{\mr{पाय}}, \textquotedblleft{}a foot\textquotedblright{}, and \emph{pīṭ}, a beating, from the family of Sanskrit \emph{piṭṭayati}, \textquotedblleft{}beats\textquotedblright{} (CDIAL 8165): beating the road with your feet.
\end{cousinweb}

\begin{grammarlens}[title={What you've built}]
One more word for getting around.
\end{grammarlens}

\subsection*{Your turn}
\begin{itemize}
\raggedright
  \item \emph{Say it:} \emph{pāypīṭ}
  \item \emph{Say it:} \emph{pāypīṭ}, once more
  \item \emph{Say it:} \emph{bas navhatī mhaṇūn āmcī khūp pāypīṭ jhālī.} — There was no bus, so we had a long trudge on foot.
  \item \emph{Recall:} say the Marathi for a brake, then the Marathi for a signpost, then say \emph{pāypīṭ} again
\end{itemize}

\begin{tcolorbox}[breakable,colback=teal!4,colframe=teal!35!black,title={Before you move on}]
What does \mr{पायपीट} mean? (\textquotedblleft{}A long trudge; a long walk\textquotedblright{}.) Say it once more.
\end{tcolorbox}
